% GENERATED FILE. Edit canonical lessons, then run npm run generate:books.
% canonical-source-hash: fnv1a64:4dc6786a34c28eca
% canonical-lessons: FR-C374-but, FR-C374-intention, FR-C374-objectif, FR-C374-resolution, FR-C374-viser

\chapter{Mes objectifs}
\label{ch:fr-mes-objectifs}
\hlchaptermodality{374}

\begin{chapteropening}
\textbf{By the end of this chapter:} \emph{I can say what I aim to do: my goal, my intention, an objective, a resolution and what I am aiming at.}

Complete the last lesson of chapter 374: I can say what I aim to do: my goal, my intention, an objective, a resolution and what I am aiming at.
\end{chapteropening}

\section[le but]{le but — the aim, the goal}

\label{lesson:FR-C374-but}



\begin{quote}
Before the new one: say the French for flexible, then the French for to consider.
\end{quote}

\subsection*{You'll want to know: le but}
\textbf{le but} — \textquotedblleft{}the aim, the goal\textquotedblright{}.

\begin{quote}
Mon but, c'est de trouver un travail à l'étranger.
\end{quote}

\begin{quote}
\emph{My aim is to find a job abroad.}
\end{quote}

In football, \emph{un but} is a goal. The final \emph{t} is often said.

\begin{cousinweb}
\emph{but} comes from Old French \emph{but} \textquotedblleft{}end, aim\textquotedblright{} (EtymDB), linked to \emph{butte} \textquotedblleft{}a mound used as a target\textquotedblright{}; it is usually traced to a Frankish word for a stump or a log.
\end{cousinweb}

\begin{grammarlens}[title={What you've built}]
One more word for telling, explaining and describing at B1.
\end{grammarlens}

\subsection*{Your turn}
\begin{itemize}
\raggedright
  \item \emph{Say it:} \emph{le but}
  \item \emph{Say it:} \emph{le but}, once more
  \item \emph{Say it:} a sentence of your own with \emph{le but}
  \item \emph{Recall:} say the French for flexible, then the French for to consider, then say \emph{le but} again
\end{itemize}

\begin{tcolorbox}[breakable,colback=teal!4,colframe=teal!35!black,title={Before you move on}]
What does le but mean? (\textquotedblleft{}The aim, the goal\textquotedblright{}.) Say it once more.
\end{tcolorbox}
\section[l'intention]{l'intention — the intention}

\label{lesson:FR-C374-intention}



\begin{quote}
Before the new one: say the French for to consider, then the French for the aim.
\end{quote}

\subsection*{You'll want to know: l'intention}
\textbf{l'intention} — \textquotedblleft{}the intention\textquotedblright{}.

\begin{quote}
Je n'ai pas l'intention de partir.
\end{quote}

\begin{quote}
\emph{I have no intention of leaving.}
\end{quote}

\emph{Avoir l'intention de} + infinitive. Feminine.

\begin{cousinweb}
\emph{intention} comes through Old French from Latin \emph{intentiō} \textquotedblleft{}a stretching, an aim\textquotedblright{} (EtymDB), from \emph{intendere} \textquotedblleft{}to stretch toward\textquotedblright{}, built on \emph{tendere} \textquotedblleft{}to stretch\textquotedblright{} (Lewis \& Short).
\end{cousinweb}

\begin{grammarlens}[title={What you've built}]
One more word for telling, explaining and describing at B1.
\end{grammarlens}

\subsection*{Your turn}
\begin{itemize}
\raggedright
  \item \emph{Say it:} \emph{l'intention}
  \item \emph{Say it:} \emph{l'intention}, once more
  \item \emph{Say it:} a sentence of your own with \emph{l'intention}
  \item \emph{Recall:} say the French for to consider, then the French for the aim, then say \emph{l'intention} again
\end{itemize}

\begin{tcolorbox}[breakable,colback=teal!4,colframe=teal!35!black,title={Before you move on}]
What does l'intention mean? (\textquotedblleft{}The intention\textquotedblright{}.) Say it once more.
\end{tcolorbox}
\section[l'objectif]{l'objectif — the objective, the target; the lens}

\label{lesson:FR-C374-objectif}



\begin{quote}
Before the new one: say the French for the aim, then the French for the intention.
\end{quote}

\subsection*{You'll want to know: l'objectif}
\textbf{l'objectif} — \textquotedblleft{}the objective, the target; the lens\textquotedblright{}.

\begin{quote}
Notre objectif est de finir avant juin.
\end{quote}

\begin{quote}
\emph{Our objective is to finish before June.}
\end{quote}

On a camera, \emph{l'objectif} is the lens. As an adjective, \emph{objectif} means objective. Masculine.

\begin{cousinweb}
\emph{objectif} comes from Medieval Latin \emph{obiectīvus} (EtymDB), from \emph{obicere} \textquotedblleft{}to throw before, to put before\textquotedblright{} (Lewis \& Short), built on \emph{iacere} \textquotedblleft{}to throw\textquotedblright{}.
\end{cousinweb}

\begin{grammarlens}[title={What you've built}]
One more word for telling, explaining and describing at B1.
\end{grammarlens}

\subsection*{Your turn}
\begin{itemize}
\raggedright
  \item \emph{Say it:} \emph{l'objectif}
  \item \emph{Say it:} \emph{l'objectif}, once more
  \item \emph{Say it:} a sentence of your own with \emph{l'objectif}
  \item \emph{Recall:} say the French for the aim, then the French for the intention, then say \emph{l'objectif} again
\end{itemize}

\begin{tcolorbox}[breakable,colback=teal!4,colframe=teal!35!black,title={Before you move on}]
What does l'objectif mean? (\textquotedblleft{}The objective, the target; the lens\textquotedblright{}.) Say it once more.
\end{tcolorbox}
\section[la résolution]{la résolution — the resolution, the firm decision}

\label{lesson:FR-C374-resolution}



\begin{quote}
Before the new one: say the French for the intention, then the French for the objective.
\end{quote}

\subsection*{You'll want to know: la résolution}
\textbf{la résolution} — \textquotedblleft{}the resolution, the firm decision\textquotedblright{}.

\begin{quote}
Ma bonne résolution pour cette année : faire du sport.
\end{quote}

\begin{quote}
\emph{My resolution for this year: to do some sport.}
\end{quote}

\emph{Les bonnes résolutions} — New Year resolutions.

\begin{cousinweb}
\emph{résolution} is a learned borrowing from Latin \emph{resolūtiō} \textquotedblleft{}a loosening, a solution\textquotedblright{} (EtymDB), from \emph{resolvere} \textquotedblleft{}to loosen\textquotedblright{}, built on \emph{solvere} \textquotedblleft{}to loosen\textquotedblright{} (Lewis \& Short).
\end{cousinweb}

\begin{grammarlens}[title={What you've built}]
One more word for telling, explaining and describing at B1.
\end{grammarlens}

\subsection*{Your turn}
\begin{itemize}
\raggedright
  \item \emph{Say it:} \emph{la résolution}
  \item \emph{Say it:} \emph{la résolution}, once more
  \item \emph{Say it:} a sentence of your own with \emph{la résolution}
  \item \emph{Recall:} say the French for the intention, then the French for the objective, then say \emph{la résolution} again
\end{itemize}

\begin{tcolorbox}[breakable,colback=teal!4,colframe=teal!35!black,title={Before you move on}]
What does la résolution mean? (\textquotedblleft{}The resolution, the firm decision\textquotedblright{}.) Say it once more.
\end{tcolorbox}
\section[viser]{viser — to aim (at)}

\label{lesson:FR-C374-viser}



\begin{quote}
Before the new one: say the French for the objective, then the French for the resolution.
\end{quote}

\subsection*{You'll want to know: viser}
\textbf{viser} — \textquotedblleft{}to aim (at)\textquotedblright{}.

\begin{quote}
Elle vise la première place.
\end{quote}

\begin{quote}
\emph{She is aiming for first place.}
\end{quote}

\emph{Viser à} + infinitive — to aim to.

\begin{cousinweb}
\emph{viser} is inherited through Old French from a Vulgar Latin \emph{vīsāre} (EtymDB), from Latin \emph{vīsere} \textquotedblleft{}to view closely\textquotedblright{} (Lewis \& Short), built on \emph{vidēre} \textquotedblleft{}to see\textquotedblright{}.
\end{cousinweb}

\begin{grammarlens}[title={What you've built}]
One more word for telling, explaining and describing at B1.
\end{grammarlens}

\subsection*{Your turn}
\begin{itemize}
\raggedright
  \item \emph{Say it:} \emph{viser}
  \item \emph{Say it:} \emph{viser}, once more
  \item \emph{Say it:} a sentence of your own with \emph{viser}
  \item \emph{Recall:} say the French for the objective, then the French for the resolution, then say \emph{viser} again
\end{itemize}

\begin{tcolorbox}[breakable,colback=teal!4,colframe=teal!35!black,title={Before you move on}]
What does viser mean? (\textquotedblleft{}To aim (at)\textquotedblright{}.) Say it once more.
\end{tcolorbox}
